\scanpage{101}
\mspage{94}

\deleted{Il est \unsure{peu|pas} \unsure{évident|immédiat} que}

On définit un système inductif sur $\widetilde C$, à valeurs dans $\underline{\mathrm{Filt}}^{\pi}$, en posant

\[
 Z_{\deleted{\unsure{\alpha}} C'_i}=\underset{\alpha\in C'_i(\unsure{\subset|\in} C_i)=\underline M^\pi/X_i)}{\varinjlim} Z_\alpha
\]

(où un $\alpha\in C_i$ est formé, je le rappelle, de $Z_\alpha\xrightarrow{u_\alpha}X_i$, avec $Z_\alpha\in M^\pi\subset\underline{\mathrm{Filt}}^\pi$).

Les $Z_{C'_i}$ ne sont plus forcément dans $M^\pi$, mais seulement dans $\mathrm{Filt}^\pi$. \deleted{\unsure{Utilisons|Utilison}} \deleted{\unsure{ainsi}} \deleted{(Alors à quoi ça a servi de \deleted{\unsure{prendre}} \inserted{définir} les $C_i$ eux-$=$s comme $\underline M^\pi/X_i$, au lieu de $\underline{\mathrm{Filt}}^\pi/X_i$ ? )}

\deleted{\unsure{des}} \unsure{Le} système inductif

\[
 (Z_{C'_i})_{C'_i\in\widetilde C_i}
\]

est \deleted{\unsure{bien}} \inserted{\unsure{parfait\textsuperscript{mt}}} adapté à la filtration \unsure{cardinale} pour le cardinal $\pi$, et il est immédiat que sa $\varinjlim$ est \inserted{\unsure{isom à|idem}} $\underset{\alpha\in\deleted{\unsure{D}} C_i=\underline M^{\pi\deleted{\unsure{i}}}/X_i}{\varinjlim} Z_\alpha=X_{\unsure{i|i'}}\in\deleted{\unsure{\varnothing}} M$.

Donc par b), l'ens \deleted{des $\unsure{\alpha}$} $\widetilde C_M$ des $C'_{\deleted{\unsure{i}} i}\in\widetilde C_i$ tels que $Z_{C'_i}\in M$ (donc $Z_{C'_i}\in M^\pi$) est cofinal dans $\unsure{\underline{\widetilde C}|\widetilde C|\widetilde C_i}$. Il s'ensuit donc que l'ens. $\widetilde C_M$ est lui aussi cofinal dans $\widetilde C$.

% archive sheet number: 100
